\documentclass[11pt]{article}
\usepackage[margin=1in]{geometry}
\usepackage{amsmath,amssymb,amsthm}
\newtheorem{theorem}{Theorem}
\newcommand{\C}{\mathcal C_\alpha}
\begin{document}
\title{A compulsory relative outlier in bounded-degree unweighted graphs}
\author{Independent unweighted extension session}\date{November 2026}
\maketitle

All inverse powers and determinants below are restricted to the common
support of connected graph Laplacians. Let $0<\alpha<1$, let $G$ be a
connected simple unweighted graph of maximum degree at most $d$, and
double an existing edge with incidence vector $b$. Write
\[
 L'=L+bb^T,\quad \tau=b^TL^\dagger b\le1,\quad
 \C=(L^\dagger)^{\alpha/2}(L'^\alpha-L^\alpha)
                                      (L^\dagger)^{\alpha/2}.
\]

\begin{theorem}
Every such update satisfies
\[
 \lambda_1(\C)\ge \frac{2\alpha}{(1+\tau)\|L\|}
                     \ge \frac{\alpha}{2d}.
\]
Consequently arbitrary determinant-admissible spectral closure is
impossible at fixed maximum degree. For $0<\varepsilon<\alpha/(2d)$,
if $m=\#\{j:\lambda_j(\C)>\varepsilon\}$, then
\[
 (m-1)\log(1+\varepsilon)
  \le \alpha\log(1+\tau)-\log\left(1+\frac{\alpha}{2d}\right).
\]
The weak inequality suffices for the bound; it is strict when $m>1$.
In particular the full determinant constant $\alpha\log2$ from arbitrary
weighted paths cannot be the small-tolerance leading count constant for
this bounded-degree class.
\end{theorem}

\begin{proof}
Put $A=L|_{\mathbf1^\perp}$, $u=b$ on this support, and $R=\|A\|$.
The standard fractional resolvent formula gives, for
$D=(A+uu^T)^\alpha-A^\alpha$ and
$m(t)=u^T(A+tI)^{-1}u$,
\[
 u^TDu=\frac{\sin(\pi\alpha)}{\pi}
                 \int_0^\infty\frac{t^\alpha m(t)^2}{1+m(t)}\,dt.
\]
Here $0\le m(t)\le\tau$ and, since $A\preceq RI$ and $\|u\|^2=2$,
$m(t)\ge 2/(R+t)$. The scalar beta integral is
\[
 \frac{\sin(\pi\alpha)}{\pi}
       \int_0^\infty\frac{t^\alpha}{(R+t)^2}\,dt
                              =\alpha R^{\alpha-1}.
\]
Thus $u^TDu\ge4\alpha R^{\alpha-1}/(1+\tau)$, whereas
$u^TA^\alpha u\le2R^\alpha$. The relative Rayleigh quotient at the
nonzero vector $A^{\alpha/2}u$ proves the first displayed bound.
An existing unit edge has effective resistance at most one, and
$R\le2d$, giving the second bound.

The determinant identity
$\sum_j\log(1+\lambda_j(\C))=\alpha\log(1+\tau)$
then proves the count statement. In particular a sufficiently long
equal-entry list with product $2^\alpha$ has all its correction entries
below $\alpha/(2d)$, so it cannot be approximated in operator norm by
this class, even after adding zeros.
\end{proof}
\end{document}
